\documentclass[11pt,a4paper]{article}
\usepackage[margin=1in]{geometry}
\usepackage[T1]{fontenc}
\usepackage[utf8]{inputenc}
\usepackage{lmodern}
\usepackage{xurl}
\usepackage{hyperref}
\usepackage{booktabs}
\usepackage{tabularx}
\usepackage{enumitem}
\usepackage{setspace}

\hypersetup{
  colorlinks=true,
  linkcolor=black,
  urlcolor=blue
}

\setstretch{1.03}
\setlist[itemize]{leftmargin=1.2cm}
\setlength{\emergencystretch}{4em}
\sloppy

\title{Weekly Coding and Implementation Progress Report\\First Six Weeks of Development}
\author{}
\date{\today}

\begin{document}
\maketitle

\section*{Overview}

This report summarizes the first six weeks of development on the MFJA
third-floor Gazebo project. During this period, the work evolved from basic
repository setup and map reconstruction into a structured simulation platform
that includes the environment model, multi-robot architecture, and the first
working shuttle implementation.

The main progression across these six weeks can be described in three stages:

\begin{itemize}
\item building and correcting the simulation environment;
\item organizing the repository and multi-robot runtime architecture;
\item implementing the first practical shuttle and switch behavior.
\end{itemize}

\section*{Six-Week Summary}

\begin{tabularx}{\textwidth}{>{\raggedright\arraybackslash}p{1.2cm} >{\raggedright\arraybackslash}p{3.4cm} X}
\toprule
\textbf{Week} & \textbf{Main Focus} & \textbf{Main Result} \\
\midrule
1 & Repository setup and map reconstruction & Established the base workflow by testing the main repository, studying multi-robot references, and restarting the CAD-based reconstruction of the MFJA map. \\
\midrule
2 & Gazebo environment usability & Integrated the floor mesh and corrected practical environment issues such as mirrored geometry, doors, room details, windows, and STL orientation. \\
\midrule
3 & Codebase organization and collaboration workflow & Worked across branches and pull requests, and improved instructions and documentation to make the changes easier to test and review. \\
\midrule
4 & Conveyor and rail environment integration & Added more real lab content, exported static parts into Gazebo-ready meshes, updated the map accuracy, and pushed a more realistic version to the main branch. \\
\midrule
5 & ROS--Gazebo software architecture for robots & Built the launch and spawning structure for world loading, robot filtering, robot state publishing, bridge generation, and namespaced robot support. \\
\midrule
6 & Shuttle and switch implementation & Explored dynamic shuttle behavior, identified its limitations, and replaced it with a first useful kinematic shuttle implementation based on spline paths and position control. \\
\bottomrule
\end{tabularx}

\section*{Week 1}

During Week 1, I began by setting up and testing the main simulation repository
for the MFJA third-floor project. I switched to the new
\texttt{mfja\_3rd\_floor\_gz} repository and reviewed the
\texttt{tb3\_multi\_robot} project as a reference for multi-robot spawning in
Gazebo Harmonic. This helped me understand how launch files build robot
creation commands and how spawning workflows are structured in practice.

In parallel, I started rebuilding the MFJA map in Autodesk Inventor. This
included adding the missing room and preparing the CAD-to-Gazebo workflow so
the environment could later be transferred into the simulation in a controlled
and reusable way.

\section*{Week 2}

During Week 2, I focused on making the simulation environment usable in
practice, not only visually correct. I exported the floor mesh as STL,
integrated it into the Gazebo model structure, and worked on the key model
files such as \texttt{mfja\_3rd\_floor.sdf} and \texttt{model.config}.

I also fixed several concrete issues in the environment:

\begin{itemize}
\item mirrored geometry;
\item incorrect door dimensions;
\item missing room details;
\item interior window integration for Room 315;
\item STL orientation problems.
\end{itemize}

This week was important because it turned the environment into something that
could actually support later robotic testing rather than remaining only a rough
visual mockup.

\section*{Week 3}

During Week 3, the work moved from only editing the environment to organizing
the codebase and development workflow itself. I worked on several branches,
opened pull requests, and started preparing my changes to be merged into the
AIP-Primeca repository main branch.

At this stage, I also started writing clearer test instructions and
documentation so that the changes, especially those related to multi-robot
support, could be understood, tested, and reviewed more easily by others. This
made the project more maintainable and more collaborative.

\section*{Week 4}

During Week 4, I integrated more real lab content into the simulation. I
received the conveyor and rail parts, identified which components were static,
and exported all static parts into a single STL file for Gazebo integration.

I also updated the map to a newer version with more accurate windows and
additional room details, pushed the current version to the main branch, and
added a detailed \texttt{README}. This was the week where the environment
became much closer to the real MFJA layout and started looking like a real
digital counterpart of the physical lab.

\section*{Week 5}

During Week 5, I worked on the actual ROS--Gazebo software structure for
running the simulation with robots. I organized the launch architecture around:

\begin{itemize}
\item a world file;
\item \texttt{robots.yaml};
\item URDF files;
\item SDF robot models;
\item dynamic spawning through \texttt{ros\_gz\_sim create}.
\end{itemize}

I also worked on the runtime flow that:

\begin{itemize}
\item loads the world;
\item filters enabled robots;
\item resolves the correct model and URDF for each robot;
\item launches \texttt{robot\_state\_publisher};
\item generates bridge configurations for \texttt{ros\_gz\_bridge}.
\end{itemize}

For mobile robots, I added logic to rewrite TIAGo topics per robot instance so
the architecture can support namespaced multi-robot execution instead of only
one robot. This was a major architectural step because it transformed the
simulation from a static world into a reusable runtime system.

\section*{Week 6}

During Week 6, I focused on the hardest implementation part so far: the shuttle
and switch system. I first tried to model the shuttle dynamically, but it did
not work properly because the shuttle got stuck and the wheel/contact-point
behavior between the shuttle and the rails was not yet solved.

Instead of remaining blocked on the dynamic model, I switched to a more useful
kinematic implementation. I extracted points for the big loop and the small
loop, exported them to Excel, created spline-based paths, and used a
set-position method in code to move the shuttle along those paths.

This gave me the first working shuttle behavior, even though it still needed
refinement. The important result of this week was that shuttle motion became
visible, testable, and ready to evolve into the more structured system that was
developed later.

\section*{Overall Progress After Six Weeks}

After these first six weeks, the project had already achieved the following
foundation:

\begin{itemize}
\item a reconstructed and corrected MFJA third-floor environment;
\item a Gazebo-ready floor and room model structure;
\item clearer repository organization and documentation;
\item a practical ROS--Gazebo architecture for multi-robot simulation;
\item namespaced support for multiple robot instances;
\item the first working shuttle implementation based on kinematic motion.
\end{itemize}

In other words, these six weeks built the technical base that made the later
Room 315 kinematic shuttle system possible.

\section*{Conclusion}

The first six weeks were mainly about building the project backbone. The work
started from repository setup and environment correction, then progressed toward
software architecture and finally reached the first operational shuttle
behavior. Although the dynamic shuttle approach was not yet successful, the
decision to move to a kinematic implementation ensured steady progress and
created a working basis for the next development stages.

\end{document}
